\subsection{Full local generization groups and adjacent flags}

\begin{result}{E.32}{Proved}{Completion on a convex pair and on its crossing stabilizers}
Let $\sigma\supset\tau$ be multicurves differing by one curve,
with the indicated pinnings.  Put
\[
 G_0=\Stab(\sigma)/\Delta_\sigma,\qquad
 G_1=\Stab(\tau)/\Delta_\tau,
 \qquad K=\Stab(\sigma,\tau)/\Delta_\tau.
                                                               \tag{E.85}
\]
The stabilizer of the corresponding arrow in the completed exit
category, under the actions of both completed endpoint groups,
is the full completion $\widehat K$.  The two homomorphisms to
the endpoints are exactly the completions of the maps in (E.85).
The same assertion holds for the automorphism group of every
specified flag of successive one-curve smoothings, with all its
endpoint actions and its projections to its individual arrows.
No congruence-completion identification is used.
\end{result}
\begin{proof}
We first establish two categorical and group-theoretic facts.
A full subcategory $D$ of a finite-object category $C$ is
\emph{convex} when every factorization $x\to z\to y$ with
$x,y\in D$ has $z\in D$.  Each finite congruence on $D$
extends to $C$: retain its quotients on pairs in $D$, and
collapse every other nonempty Hom-set to a singleton, retaining
empty Hom-sets.  Composition whose endpoints lie in $D$ occurs
entirely in $D$ by convexity.  Every other composite has its
uniquely prescribed value in a singleton.  These rules preserve
identities and associativity.  Therefore
\[
                 \widehat D\longrightarrow\widehat C
       \quad\hbox{is fully faithful}.                         \tag{E.86}
\]
This is a statement about the actual finite-congruence
completion; it follows by restricting and extending its
cofinal finite quotients.

Here is the group fact, including the exactness needed for
flags.  Let $q:B\twoheadrightarrow H$ have kernel $N$, with
the full profinite topology on $N$ induced from $B$.  For any
homomorphism $f:A\to H$, set $P=A\times_H B$.  Then
\[
 \widehat P\xrightarrow{\sim}
       \widehat A\times_{\widehat H}\widehat B.                \tag{E.87}
\]
To prove the necessary full topology, take a finite-index
subgroup $U\leq P$.  Choose a finite-index normal subgroup
$V\leq B$ with $V\cap N\subset U\cap N$.
Write $P_V$ for the inverse image of $V$ in $P$ and let
$W$ be the image of $(U\cap P_V)N$ in $A$.
Both $P_V$ and $W$ have finite index, and
\[
  P_V\cap\operatorname{pr}_A^{-1}(W)
     =(U\cap P_V)(V\cap N)\subset U.
\]
Thus $P$ has its full topology induced from $A\times B$.
Completion of $N\to P\to A$ is exact, with its first map
injective because $N$ already has its full topology in $B$.
The same is true of $N\to B\to H$.
The natural map in (E.87) is injective by the topology just
proved.  It is surjective because its map on the quotients
$\widehat A$ is onto and its map on their kernels is the
identity of $\widehat N$.  This proves (E.87).
The same assertion applies when the image of $q$ has finite
index in $H$: first replace $H$ by that image and $A$ by
its inverse image, both with their induced full topologies.

Return to the two strata.  In the skeletal category a
noninvertible arrow strictly decreases curve number, so
these two adjacent objects form a convex full subcategory.
Its crossing set is a finite union of transitive
$(G_1\times G_0)$-sets.  The finite choices are the
contraction maps of their finite labelled dual graphs.
For a chosen incidence its stabilizer is $K$ in (E.85),
with the target inclusion and the source quotient of A.6.
The latter map has image
\[
 H=\Stab(\sigma,\tau)/\Delta_\sigma\leq G_0,
 \quad [G_0:H]<\infty,
\]
and kernel $N=\Delta_\sigma/\Delta_\tau\cong\Z$.
Cutting along $\tau$ identifies a finite-index pinned
subgroup of $G_1$ with the product of the component mapping
class groups.  The residual curve is essential on its
component: a peripheral residual curve would be inessential
on the original surface or parallel to a curve of $\tau$.
The multitwist topology calculation used in E.29 applies to
that component.  Hence $N$ has its full topology in $G_1$.

Apply the diagonal-topology argument of E.29 to
$K\leq G_1$, its kernel $N$, and its quotient $H$.
It gives the embedding
\[
          \widehat K\lhook\joinrel\longrightarrow
                    \widehat G_1\times\widehat G_0.           \tag{E.88}
\]
Every finite equivariant quotient of the chosen transitive
crossing set factors through a finite quotient of
$G_1\times G_0$.  Conversely each finite quotient of this
product gives the finite coset set by the image of $K$;
these sets, together with the finite endpoint groups, are
finite categorical quotients of the two-object category.
Its completed crossing component is consequently
\[
       (\widehat G_1\times\widehat G_0)/\widehat K.             \tag{E.89}
\]
There are only finitely many incidence components, and
their tags can be retained in these quotients.  Thus the
argument applies to each component simultaneously.
Equation (E.86) identifies it with the same Hom-set in
the full completed exit category.  Formula (E.89) proves
the asserted stabilizer and the two endpoint maps.

Finally fix an actual flag of adjacent arrows.  Its
automorphism group in any category is the iterated fibre
product of its arrow stabilizers over the intermediate
automorphism groups.  On the discrete side, the source map
of each next arrow stabilizer has finite-index image and
kernel the rank-one twist group just used.  That kernel
has its full topology in the arrow stabilizer, since it
already has it in the target endpoint group.
Apply (E.87), with its finite-index variant, successively.
The resulting profinite fibre product is the full completion
of the discrete flag automorphism group.  On the completed
category side the same fibre product computes the actual
flag automorphism group by (E.89).  All the identifications
are the natural maps of tuples, so their projections to
arrows and endpoints agree exactly.  This proves the flag
assertion as well.
\end{proof}

Every noninvertible arrow of the completed category is a
composite of adjacent arrows; automorphisms are retained in
the endpoint groups.  Indeed each discrete noninvertible arrow has such a
factorization.  There are finitely many possible intervening
object lists.  The images of the corresponding finite
products of profinite Hom-sets under composition are compact,
and hence closed; their finite union contains the dense
discrete noninvertible Hom-set.  It is therefore the whole
noninvertible Hom-set.
This gives an exact reduction of crossing reconstruction
to the adjacent-arrow maps.  E.32 computes their full
stabilizers; recovering the maps themselves from the
specified operadic action remains the step under attack.
